%% ---------------------------------------------------------------------------
%% Cover letter for the Journal of Systems and Software submission.
%% One page. Every number quoted here is emitted by paper/analysis.py into
%% paper/generated/macros*.tex; keep them in step with paper/main.tex.
%% Build: pdflatex cover_letter.tex   ->  cover_letter.pdf
%% ---------------------------------------------------------------------------
\documentclass[11pt]{article}

\usepackage[T1]{fontenc}
\usepackage[utf8]{inputenc}
\usepackage{times}
\usepackage[margin=0.9in]{geometry}
\usepackage[hidelinks]{hyperref}
\usepackage{enumitem}
\pagestyle{empty}
\setlength{\parindent}{0pt}
\setlength{\parskip}{0.32em}
\setlist[itemize]{leftmargin=1.1em,itemsep=0pt,topsep=0.15em,parsep=0pt}

\begin{document}

{\small\textbf{Fahad Hafeez}\\
Newcastle University, Newcastle upon Tyne, NE1 7RU, United Kingdom\\
\href{mailto:f.hafeez2@newcastle.ac.uk}{f.hafeez2@newcastle.ac.uk}}

\vspace{0.2em}
\hrule
\vspace{0.3em}

19 November 2026

The Editors-in-Chief, \emph{Journal of Systems and Software}

\textbf{Re: submission of \emph{Agent Recovery Bench: Measuring Whether Failures Are
Detected Early and Repaired Reliably in Tool-Using and Multi-Agent Workflows}}

Dear Editors,

I am pleased to submit the above manuscript for consideration as a research article in the
\emph{Journal of Systems and Software}.

Software built from LLM agents is judged almost entirely by task success: did the workflow
finish? That answers none of the questions a reliability engineer has to ask---why it failed,
whether the failure was noticed, and whether the repair was worth its cost. Existing
benchmarks score what a system reports, so a plausible claim of success counts the same as a
verified one. The manuscript addresses this with \emph{Agent Recovery Bench}, contributing:
\begin{itemize}
  \item a seven-class fault taxonomy, each class tied to its identifying signal and repair
        operator;
  \item four deterministic tool environments where success is graded from world state, never
        from the agent's own report;
  \item a fault-injection harness that replays an identical reference trajectory with and
        without a scheduled fault, screening out injections that cause no harm;
  \item a protocol comparing no recovery, retry-only, rule-based and model-based recovery under
        a fixed price sheet.
\end{itemize}

Across 861 harmful fault pairs and 4,920 episodes, a rule-based detector reaches step-level AUC
0.772 (F1 51.5) and rule-based recovery repairs 60.9\% of harmful failures, against 0.0\% with
no recovery and 0.0\% with retries alone; a learned ranker reaches 51.3\% and does not beat the
canonical rule mapping---a negative result, since a learned policy is the current default.
Repair is class-dependent (about 96\% for misuse
and contradictory evidence, 41\% for planning error) and timing-dependent (71.7\% when the fault
is detected before it propagates, 50.8\% when detection is late). Agents claim success 45.6\% of
the time but succeed 28.0\% of the time, so 38.4\% of success claims are wrong.

The work fits the journal's coverage of software engineering for AI systems and of software
reliability: the contribution is not a new agent but a measurement instrument and an empirical
evaluation, with the harness, evaluator, fault schedule and analysis script released so that
later work can extend or contest every number. All environments are simulated, and the
manuscript states that it measures relative reliability under controlled faults, not production
reliability.

The manuscript is original, is not published or under consideration elsewhere, and I am its sole
author, with no competing interests and no funding. The dataset, detector, ranker and results are
openly available from the Hugging Face Hub and the harness, evaluator and analysis script from
GitHub under the MIT licence (tag \texttt{arb-1.0.1}); no credential is needed to reproduce any
reported number.

Thank you for considering the submission; correspondence should be addressed to me at the e-mail
above.

Yours sincerely,

\vspace{1.2em}
Fahad Hafeez\\
Newcastle University

\end{document}
